% DRAFT (scratch) for 06_theory.tex, after the "never fires" paragraph, before
% "Which Assumptions Can Be Dropped".  Numbers wait on prechecks/timeout_rule.

\subsection{A Clock Keeps the Same Promise}
\label{sec:clock}

The rule deployed systems use against starvation is a maximum waiting time. Once the
oldest waiting job has waited $\theta$, it is served next; until then the base policy
decides. Call this Timeout($\theta$). It reads arrival times and a clock, and never a
service time. It turns out to carry the same kind of promise as the guard.

\begin{proposition}[Timeout bound]
\label{prop:clock}
For every base policy, every input and every job $i$ with at least one overtaker,
Timeout($\theta$) gives $\mathrm{In}_i < k(\theta + L)$, and every job waits
\[
  W \;<\; W_{\mathrm{FCFS}}[i] + \theta + (3 - 2/k)L .
\]
\end{proposition}

\begin{proof}[Proof idea]
From $a_i + \theta$ on, the oldest waiting job has waited at least $\theta$, so every
dispatch until $s_i$ serves a job older than $i$. Each overtaker of $i$ therefore starts
before $a_i + \theta$ and ends less than $L$ later. The overtakers that share a server
run one after another inside $[a_i, a_i + \theta + L)$, and summing over the $k$
servers gives the bound on $\mathrm{In}_i$. Corollary~\ref{cor:sufficient} turns it
into the bound on the wait. Supplementary Section~S4.12 gives the proof in full.
\end{proof}

At $\theta = G - (3-2/k)L$ the clock makes the promise that the guard makes at
$B_{\max} = k\theta$. The two rules differ in what they charge a waiting job for. The
clock charges it for all the work the servers do while it waits, $k(t - a_q)$. The
guard charges it only for completed work of jobs that arrived after it. The
difference is the work of older jobs, which first-come first-served would have done
anyway. Near saturation that work is large. A job that arrives behind a backlog is
charged for the backlog, the clock fires for most of the queue, and the schedule falls
back to arrival order at the loads where ordering by cost pays most.
Section~\ref{sec:exp_guard} measures the difference.

Firing on either rule keeps both promises, because each proof reads only its own
condition. The clock's half needs no completion report, so a guard that also fires at
age $\theta$ keeps $W < W_{\mathrm{FCFS}} + \theta + (3-2/k)L$ when reports are late
or lost.

\subsection{An Absolute Bound from the Arrival Log}
\label{sec:absolute}

A promise relative to first-come first-served is one a student cannot check, since
nobody runs the counterfactual queue. The arrival log alone gives an absolute one. Let
$V_i^-$ be the work that a single server of rate $k$, fed the same jobs, would still hold
at $a_i$ from the jobs of rank below $i$. It follows the recursion
$V_i^- = \max\bigl(0,\ V_{i-1}^- + C_{i-1} - k(a_i - a_{i-1})\bigr)$, and its largest
value plus the arriving job, $\sigma = \max_i (V_i^- + C_i)$, is the depth of the
smallest token bucket of rate $k$ that the arrivals, counted in work, fit. That is the
burst parameter of network calculus~\cite{cruz1991calculus,leboudec2001network}.

\begin{proposition}[Absolute bound]
\label{prop:absolute}
For every non-preemptive work-conserving policy and every job $i$,
$k\,W[i] \le V_i^- + (k-1)L + \mathrm{In}_i$. Hence
$W_{\mathrm{FCFS}}[i] \le (V_i^- + (k-1)L)/k$, and under Algorithm~\ref{alg:guard}
with cap $B_{\max}$
\[
  W \;<\; \frac{V_i^- + B_{\max}}{k} + \Bigl(2 - \frac{1}{k}\Bigr)L
    \;\le\; \frac{\sigma + B_{\max}}{k} + \Bigl(2 - \frac{1}{k}\Bigr)L .
\]
\end{proposition}

\begin{proof}[Proof idea]
The busy-period accounting in the proof of Theorem~\ref{thm:identity} gives
$kW[i] \le R_i + \mathrm{In}_i$. A work-conserving $k$-server queue never holds more
than $(k-1)L$ above the fluid server, because whenever it does, all $k$ servers are
busy and it drains at the fluid rate; so $R_i \le V_i^- + (k-1)L$.
Theorem~\ref{thm:guard} bounds $\mathrm{In}_i$. Supplementary Section~S4.12 gives the
proof in full.
\end{proof}

At $k = 1$ the fluid server is the queue, $V_i^- = W_{\mathrm{FCFS}}[i]$, and the
proposition is Theorem~\ref{thm:guard} again. For $k \ge 2$ it is sharper than going
through $W_{\mathrm{FCFS}}$, which would pay the $2(k-1)L$ of
Theorem~\ref{thm:identity} a second time. An operator who promises that no job waits
longer than $D$ in a period whose arrivals fit a bucket of depth $\sigma$ sets
$B_{\max} = kD - \sigma - (2k-1)L$: the capacity over the deadline, less the burst,
less $2k-1$ job-limits. The clock obeys the same bound with $B_{\max}$ replaced by
$k\theta$. Section~\ref{sec:exp_bound} measures how close the bound comes on the traces
and how well one period's $\sigma$ predicts the next.
